\documentclass[12pt]{article}
\usepackage[utf8]{inputenc}
\usepackage[T1]{fontenc}
\usepackage{hyperref}
\title{Zhoa Language Resources for K-12 Curriculum Development}
\author{Compiled by Claude Sonnet 5.5 (claude-sonnet-5-5) for Andr\'as Kornai}
\date{}
\begin{document}
\maketitle

\section{Introduction}
Zhoa (ISO 639-3 zhw, Glottocode zhoa1238) is a Grassfields Bantu language (Ring group, Bantoid) spoken in Fungom Subdivision, Menchum Division, North-West Region of Cameroon. Wikipedia gives about 2,000 speakers (1995 figure). Documentation is very thin. It is not known to be taught in schools.

\section{Primary Educational Resources}
No primers, textbooks or school materials were found in search results.

\section{Linguistic and Reference Materials}
\begin{itemize}
\item Wikipedia article on the Zhoa language: \url{https://en.wikipedia.org/wiki/Zhoa_language}
\item Wikipedia article on the Zhoa people and town: \url{https://en.wikipedia.org/wiki/Zhoa}
\item UNESCO World Atlas of Languages entry, Zhoa in Cameroon: \url{https://en.wal.unesco.org/countries/cameroon/languages/zhoa}
\end{itemize}
No grammar or dictionary was found.

\section{Religious and Community Resources}
Joshua Project lists the Zhoa people of Cameroon, with no known Scripture translation and no published resources: \url{https://joshuaproject.net/people_groups/18991}

\section{Digital and Audio Resources}
No audio, corpus or software resources were found.

\section{Curriculum Development Context}
Almost nothing exists for Zhoa. Curriculum work would have to begin with documentation (a word list, an orthography and recordings), done together with the community.

\end{document}